\section{Experiments}
\label{sec:main-experiments}

This section provides empirical demonstration of the proposed theories. Appendix~\ref{sec:appendix-experiments} gives the simulation details. We estimate the average treatment effect $\tau$ through PROBE (Algo.~\ref{alg:probe-search}) and report the mean absolute error (MAE) over independently generated datasets. We compare PROBE with two groups of methods. The first group needs no proxy roles: adjustment for $X$, adjustment for $(X,W)$, CEVAE \citep{louizos2017causal}, and, for image proxies, adjustment for a representation that a convolutional network learns from all pixels, without holding out a block or checking balance. The second group needs the designated roles for each proxy: proximal two-stage least squares (2SLS) \citep{tchetgen2024ProximalCausalInference}, which requires designating each proxy as treatment-side or outcome-side (Fig.~\ref{fig:causal-structures}(c)); single-proxy control \citep{park2024single}, which requires designating one proxy of the untreated outcome $Y(0)$; and its kernel version KSPC \citep{xu2025kernel}, which requires the same designation. We run each of these methods once with the correct roles and once with wrong roles. An oracle that adjusts for $(X,U)$, including the hidden confounder, serves as a reference.

\paragraph{Synthetic Data Analysis.} In our example scenario, a hidden severity $U$ affects both the treatment and the outcome, and the true effect is $\tau=1$. The proxy has five blocks. Blocks $W_1$, $W_2$, and $W_3$ are noisy measurements of $U$ and are valid held-out blocks: each satisfies Assumptions~\ref{ass:joint-latent-exchangeability} and~\ref{ass:outcome-relevant-completeness}. Blocks $W_4$ and $W_5$ are not valid: $W_5$ violates held-out treatment independence (Assumption~\ref{ass:joint-latent-exchangeability}), so no representation can balance it, and $W_4$ can be balanced but violates outcome-relevant completeness (Assumption~\ref{ass:outcome-relevant-completeness}), so it points to a wrong effect. PROBE is not told which blocks are valid, whereas the role-requiring methods are given the valid blocks in their correct-role runs. We render this structure with $(X,W)$ of dimension $8$, $88$, and $1{,}312$ (SCM-1 to SCM-3) and with six Shapes3D images in place of the six proxy coordinates of SCM-1 (SCM-4), and generate $100$ data sets of size $n=24{,}000$ for each. 

\begin{figure}[t]
\centering
\includegraphics[width=\linewidth]{probe-experiments-v5.pdf}
\caption[Errors in the synthetic and semi-synthetic data.]{Estimate minus the true effect on a symmetric log scale; each dot is one data set. Filled and hollow dots show methods that need proxy roles with correct and wrong roles; on the semi-synthetic data, single-proxy control is used as an ATE estimator.}
\label{fig:experiments}\label{fig:synthetic}\label{fig:realdata}
\end{figure}
Fig.~\ref{fig:synthetic} (left) shows that PROBE, without being told the roles, stays near zero error in all four settings, while adjustment for $X$ or $(X,W)$ stays biased and CEVAE degrades as the dimension grows. Methods that need the roles are accurate only with the correct roles and fail badly with wrong ones. The invalid block $W_4$ passes the balance screen in $396$ of $400$ data sets, yet plurality aggregation over held-out blocks removes it in every one of them, as Thm.~\ref{thm:unknown-valid-block-identification} predicts.

\paragraph{Real-World Data Analysis.}\label{sec:real-world-analysis} We use four real-world data sets. Twins \citep{louizos2017causal}, IHDP \citep{hill2011bayesian}, and ACIC 2016 \citep{dorie2019automated} are semi-synthetic benchmarks built on real covariates, with $11{,}984$, $747$, and $4{,}802$ units and $42$, $24$, and $78$ covariates. In each, we hide one covariate as $U$ and generate the proxies and the treatment from it with one fixed recipe, so the effect is known exactly. The right heart catheterization (RHC) study \citep{connors1996effectiveness} is observational, with $5{,}735$ critically ill patients, days survived within $30$ as the outcome, and ten physiological measurements grouped into five proxy blocks, as analyzed by \citet{cui2024semiparametric}. No true effect is known for the RHC study.

Fig.~\ref{fig:realdata} (right) shows that PROBE, without being told the roles, is the most accurate method on Twins and among the three most accurate on IHDP and ACIC. As in the synthetic data, methods that need the roles are accurate only with the correct roles and degrade sharply with wrong ones. On RHC, where no true effect is known, PROBE's interval contains the proximal estimate that \citet{cui2024semiparametric} obtain by assigning the roles, while KSPC's estimate changes with the block it is given~(Fig.~\ref{fig:appendix-realdata}).

Overall, PROBE is accurate without being told the proxy roles, and it avoids the large failures of role-requiring methods when the roles are wrong. Together, these results indicate that proximal balancing, through PROBE, gives accurate effect estimates in high-dimensional proximal settings without requiring the proxy roles. 
